% SHARD-ID: LB-09-TWO-MAGNON
% SHARD-TITLE: The two-magnon sector: Bethe oracle and the exact soft slope
% SHARD-SOURCES: definitions.md D6--D8; theory/oracle-bethe.md;
%   theory/soft-current-recon.md; theory/ml2-completeness.md;
%   theory/spin-s-twomagnon.md; claims/CLAIMS.md rows S2, S2-2body,
%   S2-2body-S, OR1, OR2, N1, ML2
% SHARD-CLAIMS: D6, D7, D8, OR2, ML2, S2, S2-2body, S2-2body-S, OR1, N1

\section{The two-magnon sector: Bethe oracle and the exact soft slope}
\label{sec:two-magnon}

Everything else in this labbook is an attempt to say something general. This
section is the opposite: it is the one place where the campaign has a model it
can solve exactly, and where a proposed soft law can therefore be checked
against arithmetic rather than against intuition. The model is the isotropic
Heisenberg ferromagnet on a chain. Its fully polarised state is an exact
ground state, its one-magnon sector is a free hopping problem, and its
two-magnon sector reduces --- fibre by fibre in the total momentum --- to a
half-line Jacobi matrix with a single boundary defect. Nothing in this section
needs integrability as a hypothesis: every statement below is obtained by
writing down the lattice eigenvalue equations and solving them.

Three features of this referee should be kept in mind while reading.

First, the magnon dispersion is quadratic at small momentum,
$\omega(k)=J(1-\cos k)$, so ``soft'' here means small \emph{momentum}, and the
energy variable is a bad coordinate: the phase is analytic in $k_s$ and only
Puiseux (half-integer powers of $\omega_s$) in the energy. Statements are
therefore always organised in momentum.

Second, what is exactly true in the soft limit is weaker than the field-theory
slogan. The two-magnon scattering ratio tends to one as the soft momentum
vanishes --- Dyson's decoupling of a zero-momentum magnon --- but this is a
plain limit, not an Adler zero: no theorem here says that a soft insertion is
annihilated, only that the connected amplitude becomes the identity, with a
computable first-order departure.

Third, that first-order departure is the campaign's actual number. On the
physical outgoing/incoming channel the phase behaves as
$\delta_{\mathrm{phys}}=2\,\sgn(v_h-v_s)\,k_s+O(k_s^2)$ for spin $1/2$, and,
in the exact generalisation to every site spin $S$, as
$\sgn(v_h-v_s)\,k_s/S+O(k_s^2)$. The slope $1/S$ is the sharpest exactly
proved statement the campaign owns.

\subsection{The oracle model and its conventions}

The next three definitions are frozen: they are the shared vocabulary of every
later two-magnon statement, and they are reproduced here in full so that this
document can be read without opening the definitions file.

\begin{definition}[The Heisenberg-ferromagnet oracle model and its bases]
\label{def:fm-oracle-model}
On the infinite chain --- or on a periodic ring of $N$ sites, whenever finite
volume is stated --- put $\bm{S}_x=\bm{\sigma}_x/2$, take $J>0$, and let
\begin{equation}
  H=\sum_x h_{x,x+1}
   =-J\sum_x\Bigl(\bm{S}_x\cdot\bm{S}_{x+1}-\tfrac14\Bigr),
  \qquad
  h_{x,x+1}=\tfrac{J}{2}\bigl(1-P_{x,x+1}\bigr),
\end{equation}
where $P_{x,x+1}$ swaps the two spins. The vacuum is the fully polarised state
$\ket{\Omega}=\ket{\uparrow\uparrow\cdots}$. The coordinate bases are
$\ket{x}=S^-_x\ket{\Omega}$ and $\ket{x,y}=S^-_xS^-_y\ket{\Omega}$ with $x<y$.
Lattice momenta lie in $(-\pi,\pi]$, we write $z(k)=e^{ik}$, and
\begin{equation}
  \ket{k}_N=N^{-1/2}\sum_{x=0}^{N-1}e^{ikx}\ket{x},
  \qquad k=2\pi n/N .
\end{equation}
On the infinite chain the same sum denotes the delta-normalised generalised
state. The symbols $\omega(k)$ and $v(k)=d\omega/dk$ denote the one-magnon
energy and group velocity derived from this $H$.
\end{definition}
\provenance{D6}{\texttt{definitions.md} D6 (FROZEN --- not editable without
reopening the oracle review loop)}{\texttt{theory/checks/oracle\_bethe\_check.py}}%
{\texttt{theory/verdicts/oracle-bethe-r2.md}}

\begin{definition}[Ordered-coordinate Bethe wave and the scattering convention]
\label{def:bethe-convention}
In the chamber $x<y$ the two-magnon coordinate Bethe wave is
\begin{equation}
  \psi(x,y)=A_{12}\,e^{i(k_1x+k_2y)}+A_{21}\,e^{i(k_2x+k_1y)},
\end{equation}
and one sets
\begin{equation}
  S_{12}(k_1,k_2):=\frac{A_{12}}{A_{21}},
  \qquad
  S_{21}:=\frac{A_{21}}{A_{12}},
\end{equation}
while $E(k_1,k_2)$ denotes the corresponding scattering energy. For real wave
packets with $v(k_2)>v(k_1)$ and the $k_2$ packet initially on the left,
$A_{21}$ is the \emph{incoming} coefficient and $A_{12}$ the \emph{outgoing}
one; reversing the channel convention gives $S_{21}=S_{12}^{-1}$. The rapidity
is
\begin{equation}
  \lambda(k):=\tfrac12\cot(k/2).
\end{equation}
Whenever $S_{12}=e^{i\delta}$ is expanded at a point where $S_{12}=1$, the
symbol $\delta$ means the unique continuous real branch vanishing at that
point. For a two-magnon bound state we write $K:=k_1+k_2$ for the total
momentum, $r:=y-x$ for the separation, $f_r$ for the relative-coordinate wave,
$t$ for its geometric decay factor, $\eta:=-\log\abs{t}$, and $E_b(K)$ for its
energy.
\end{definition}
\provenance{D7}{\texttt{definitions.md} D7 (FROZEN)}%
{\texttt{theory/checks/oracle\_bethe\_check.py}}%
{\texttt{theory/verdicts/oracle-bethe-r2.md}}

\begin{definition}[The soft limit of the Bethe oracle]
\label{def:bethe-soft-limit}
The hard momentum $k_h\in(0,\pi)$ is held fixed while the \emph{signed} soft
momentum $k_s\to0$. Write $\omega_h:=\omega(k_h)$, $v_h:=v(k_h)$,
$\omega_s:=\omega(k_s)$, $v_s:=v(k_s)$, and $\sigma:=\sgn(k_s)$ when a
one-sided energy expansion is used. On this half-zone $v_h>v_s$ for all
sufficiently small $\abs{k_s}$, so the ratio $S_{12}$ of
\cref{def:bethe-convention} is the physical outgoing/incoming ratio there. A
simultaneous limit of $k_h$ and $k_s$, or a hard momentum outside this
half-zone, is \emph{not} included unless explicitly stated.
\end{definition}
\provenance{D8}{\texttt{definitions.md} D8, FROZEN-AS-AMENDED: the naming of
$\omega_h,v_h,v_s$ and the remark that $v_h>v_s$ for small $\abs{k_s}$ were
added in oracle round r2 and break no consumer}%
{\texttt{theory/checks/oracle\_bethe\_check.py}}%
{\texttt{theory/verdicts/oracle-bethe-r2.md} (residue 3)}

The exact solution of the model, in this vocabulary, is the following block of
facts. They are used repeatedly below and are stated once here.

\begin{oracle}[Exact one- and two-magnon data of the ferromagnet]
\label{thm:bethe-oracle-facts}
\statusProved\ Under \cref{def:fm-oracle-model,def:bethe-convention}:
\begin{align}
  \omega(k)&=J(1-\cos k), & v(k)&=J\sin k, \\
  E(k_1,k_2)&=\omega(k_1)+\omega(k_2)=J(2-\cos k_1-\cos k_2), & &
\end{align}
the exact coefficient ratio in the ordered chamber is
\begin{equation}
  S_{12}(k_1,k_2)
  =-\,\frac{e^{i(k_1+k_2)}-2e^{ik_1}+1}{e^{i(k_1+k_2)}-2e^{ik_2}+1}
  =\frac{\lambda_1-\lambda_2+i}{\lambda_1-\lambda_2-i},
  \label{eq:exact-ratio-half}
\end{equation}
with $\lambda_j=\lambda(k_j)$, and it has unit modulus for real momenta; and
there is, up to scale, exactly one two-magnon bound state per total momentum,
\begin{equation}
  f_r=t^{\,r-1},\qquad t=\cos(K/2),\qquad E_b(K)=J\sin^2(K/2),
  \label{eq:bound-band}
\end{equation}
for $0<\abs{K}\le\pi$, degenerating at $K=0$ into a non-normalisable threshold
resonance.
\end{oracle}
\provenance{oracle facts O1--O5 of \texttt{theory/oracle-bethe.md} (these have
no separate row in \texttt{claims/CLAIMS.md}; the campaign rows that consume
them are OR1 and OR2)}{\texttt{theory/oracle-bethe.md} steps (1.1)--(1.2),
collected at (1.4)}{\texttt{theory/checks/oracle\_bethe\_check.py}}%
{\texttt{theory/verdicts/oracle-bethe-r2.md}}

\emph{Idea of proof.} For two spin halves $P=2\bm{S}_1\cdot\bm{S}_2+1/2$, so
each bond term is $\tfrac{J}{2}(1-P)$, has eigenvalues $0$ and $J$, and
annihilates any parallel pair; hence $H\ket{\Omega}=0$ and $H\ge0$. Acting on
one overturned spin leaves a nearest-neighbour hopping problem, which the
Fourier states diagonalise. For two overturned spins the bond algebra gives a
free equation away from contact --- whence energy additivity --- and a single
modified equation at $y=x+1$; substituting the two-plane-wave form into that
one contact equation gives a linear relation between $A_{12}$ and $A_{21}$,
and hence \cref{eq:exact-ratio-half}. In a fixed total-momentum fibre the
relative equation is a constant-coefficient second-order recursion whose two
characteristic roots multiply to one; square summability therefore selects the
single root inside the unit disc, and matching the $r=1$ equation fixes both
the decay factor and the energy in \cref{eq:bound-band}.

\subsection{Soft decoupling: what is actually true in the limit}

\begin{oracle}[A zero-momentum magnon decouples]
\label{thm:soft-decoupling}
\statusProved\ On the half-zone of \cref{def:bethe-soft-limit} --- fixed
$k_h\in(0,\pi)$, signed $k_s\to0$ ---
\begin{equation}
  \lim_{k_s\to0^+}S_{12}(k_s,k_h)
  =\lim_{k_s\to0^-}S_{12}(k_s,k_h)=1 .
\end{equation}
\end{oracle}
\begin{scope}
This is a plain limit and nothing more. It says that the connected two-magnon
amplitude degenerates to the identity as the soft leg becomes a zero mode; it
is explicitly \emph{weaker} than an Adler-zero theorem, since it asserts no
vanishing of a soft insertion, no factorisation, and no statement about any
process other than exact two-magnon scattering in this model. The rate at
which the limit is approached --- the content of
\cref{thm:two-body-soft,thm:spin-s-slope} --- is a separate statement.
\end{scope}
\provenance{OR2 (oracle fact O6)}{\texttt{theory/oracle-bethe.md} step
(1.3).(2).2}{\texttt{theory/checks/oracle\_bethe\_check.py}; the recorded run
gives $\abs{S_{12}-1}=2.966\cdot10^{-16}$ at the sampled soft
points}{\texttt{theory/verdicts/oracle-bethe-r2.md}}

\emph{Idea of proof.} Put $k_s=0$ in \cref{eq:exact-ratio-half}. The numerator
becomes $e^{ik_h}-1$ and the denominator $1-e^{ik_h}$; the outer minus sign
turns their ratio into $1$, and the denominator is nonzero because $k_h$ is
fixed inside $(0,\pi)$. The statement is thus simply $S_{12}-1\to0$, obtained
by subtracting the identity amplitude, with no low-energy theorem involved.

\subsection{Completeness of the two-magnon sector}

The soft expansions below are statements about a distinguished channel inside
a Hilbert space. They are only meaningful once one knows what else that space
contains. That is the content of the next theorem, which is proved by direct
diagonalisation and, deliberately, without assuming that Bethe states are
complete.

\begin{theorem}[Complete two-magnon resolution, without assumed Bethe completeness]
\label{thm:two-magnon-completeness}
\statusProved\ Assume \cref{def:fm-oracle-model,def:bethe-convention,def:bethe-soft-limit},
$J>0$, and, in finite volume, $N>3$.
\begin{enumerate}
\item \emph{Finite ring.} Every nonzero two-magnon Bethe vector on the
  $N$-site ring is one of the real-pair, complex-pair (two-string), or
  $SU(2)$-descendant vectors classified in the proof. For even $N$ the
  singular completion supplies exactly one further vector, the normalised
  contact state
  \begin{equation}
    \ket{\chi_\pi}=N^{-1/2}\sum_x(-1)^x\ket{\{x,x+1\}},
    \qquad H\ket{\chi_\pi}=J\ket{\chi_\pi},
    \label{eq:singular-vector}
  \end{equation}
  of total momentum $\pi$, which carries the compactified rapidity pair
  $\{+i/2,-i/2\}$ and has no finite two-plane-wave normalisation. The formal
  coincident solutions are identically zero vectors and supply nothing. The
  physical vectors number exactly $N(N-1)/2$ and can be normalised to an
  orthonormal eigenbasis.
\item \emph{Infinite chain.} The two-magnon Hilbert space is the orthogonal
  sum of the scattering representation of \cref{def:bethe-convention} and one
  bound-magnon band. In the centre-momentum chart, with $c=\cos(K/2)$, the
  fibre is the half-line Jacobi operator
  \begin{equation}
    (h_Kf)_1=Jf_1-Jcf_2,\qquad
    (h_Kf)_r=2Jf_r-Jc(f_{r-1}+f_{r+1})\quad(r\ge2),
    \label{eq:jacobi-fibre}
  \end{equation}
  the generalised eigenfunctions and energies are
  \begin{align}
    \Psi^{\mathrm{sc}}_{Kq}(x,x+r)
      &=\frac{e^{iK(x+r/2)}}{2\pi}\bigl[S(K,q)e^{iqr}+e^{-iqr}\bigr],
      & 0&<q<\pi, \label{eq:psi-sc}\\
    S(K,q)&=\frac{c-e^{-iq}}{e^{iq}-c},
      & E(K,q)&=J(2-2c\cos q), \label{eq:S-fibre}\\
    \Psi^{\mathrm{b}}_{K}(x,x+r)
      &=\frac{e^{iK(x+r/2)}}{\sqrt{2\pi}}\sqrt{1-c^2}\;c^{\,r-1},
      & E_b(K)&=J\sin^2(K/2), \label{eq:psi-b}
  \end{align}
  and for every chamber vector $\psi$ the resolution of the identity holds:
  \begin{equation}
    \norm{\psi}^2
    =\int_{-\pi}^{\pi}\!dK\int_0^\pi\!dq\;
      \abs{\braket{\Psi^{\mathrm{sc}}_{Kq},\psi}}^2
    +\int_{-\pi}^{\pi}\!dK\;
      \abs{\braket{\Psi^{\mathrm{b}}_{K},\psi}}^2 .
    \label{eq:plancherel}
  \end{equation}
  The Plancherel measures are $dK\,dq/(2\pi)^2$ and $dK/(2\pi)$ for the
  unnormalised kernels. The spectrum on this sector is purely absolutely
  continuous; the bound state of \cref{eq:bound-band} becomes a band, present
  for every $K\ne0$ and degenerating at $K=0$ into a threshold resonance.
\item \emph{Consequence.} The finite-volume expansion of the charge-created
  state in the Bethe family is therefore \emph{unconditional}, with the
  singular class separated explicitly: for the charge $Q_{k_s}=\sum_xe^{ik_sx}S^-_x$,
  \begin{equation}
    Q_{k_s}\ket{k_h}_N
    =\sum_{B}c_B\ket{B}
     +\mathbf 1\{2\mid N,\ k_s+k_h\equiv\pi\}\;
      2i\cos\!\bigl(\tfrac{k_h-k_s}{2}\bigr)\ket{\chi_\pi},
    \label{eq:charge-expansion}
  \end{equation}
  the sum running over the normalised real, string and descendant states of
  the fibre. At $k_s=0$ the expansion collapses onto the single
  root-at-infinity descendant $Q_0\ket{k_h}_N/\sqrt{N-2}$.
\end{enumerate}
\end{theorem}
\begin{scope}
This is a generalised-eigenfunction (Plancherel) statement. It does \emph{not}
construct M\o{}ller operators and does \emph{not} prove asymptotic
completeness in the scattering-theoretic sense; that is the separate content
of \cref{thm:two-magnon-wave-operators} in
\cref{sec:wave-operators}. Nothing here is claimed for three or more magnons.
\end{scope}
\provenance{ML2}{\texttt{theory/ml2-completeness.md} steps (1.1)--(1.7); the
finite-ring count is (1.3)--(1.4), the infinite-chain resolution (1.5), the
charge-expansion consequence (1.6)}%
{\texttt{theory/checks/ml2\_completeness\_check.py}: eleven ring sizes,
maximum spectral mismatch $8.882\cdot10^{-15}$, eigenvector residual
$4.044\cdot10^{-14}$, projector error $4.640\cdot10^{-14}$, singular-overlap
error $6.280\cdot10^{-15}$}{\texttt{theory/verdicts/ml2-r2.md} (PASS, 10/10
mutants killed)}

\emph{Idea of proof.} Translation invariance fibres the two-magnon sector by
total momentum. On the ring one counts the relative basis vectors in each
fibre --- allowing for the identification of a pair with its reflection and
for the special self-opposite separation at even $N$ --- and the dimensions
sum to $N(N-1)/2$. In each fibre the Hamiltonian is a real Jacobi matrix whose
interior recursion is the free one and whose first row carries the contact
defect; its eigenvalue problem is solved through the Chebyshev-like
polynomials of the recursion, and the resulting characteristic equation is
matched, root by root, with the Bethe quantisation condition. The formal
coincident roots are shown to produce identically vanishing coordinate waves,
and at even $N$ the $K=\pi$ fibre is diagonal, producing the singular contact
vector \cref{eq:singular-vector} explicitly rather than as a limit of
two-plane-wave states. On the infinite chain the same fibration gives the
half-line Jacobi operator \cref{eq:jacobi-fibre}; a direct unit-circle
integration of the scattering waves against each other yields the fibre
completeness identity
$\delta_{rs}=\int_0^\pi\frac{dq}{2\pi}\overline{w_q(r)}w_q(s)+(1-c^2)c^{\,r+s-2}$,
in which the bound term is exactly the residue at the pole $z=c$; combining
this with Fourier--Plancherel in the centre coordinate gives
\cref{eq:plancherel}. Absolute continuity follows because the energy maps
$E(K,q)$ and $E_b(K)$ have nonvanishing gradient off a Lebesgue-null set.
Finally, Parseval in the resulting orthonormal fibre basis, together with the
explicit overlap of the charge-created state with $\ket{\chi_\pi}$, gives
\cref{eq:charge-expansion}.

\subsection{The two-body soft theorem}

We can now state the soft law of the spin-$1/2$ ferromagnet in the form the
campaign uses. The point of the derivation is that it does \emph{not} go
through the closed Bethe ratio: it uses only the local current and the single
contact equation, so the closed formula of \cref{thm:bethe-oracle-facts} is
available afterwards as an independent check rather than as an input.

\begin{theorem}[Two-body soft expansion for the spin-half ferromagnet]
\label{thm:two-body-soft}
\statusProved\ Assume \cref{def:fm-oracle-model,def:bethe-convention,def:bethe-soft-limit}
with $k_h$ fixed, $0<\abs{k_h}<\pi$, and signed $k_s\to0$. Let
$S_{\mathrm{phys}}$ be the physical outgoing/incoming ratio --- that is,
$S_{12}$ when $\sgn(v_h-v_s)=+1$ and $S_{21}=S_{12}^{-1}$ when
$\sgn(v_h-v_s)=-1$ --- and let $\delta_{\mathrm{phys}}$ be its continuous
phase vanishing at $k_s=0$. Then
\begin{align}
  \delta_{\mathrm{phys}}
    &=2\,\sgn(v_h-v_s)\,k_s+\frac{\abs{v_h}}{\omega_h}k_s^2+R_\delta,
    \label{eq:delta-phys}\\
  S_{\mathrm{phys}}
    &=1+2i\,\sgn(v_h-v_s)\,k_s
      +\Bigl[\,i\frac{\abs{v_h}}{\omega_h}-2\Bigr]k_s^2+R_S ,
    \label{eq:S-phys}
\end{align}
where the quadratic coefficient has the equivalent invariant forms
\begin{equation}
  \frac{\abs{v_h}}{\omega_h}=\frac{2J-\omega_h}{\abs{v_h}}
  =\sqrt{\frac{2J-\omega_h}{\omega_h}},
  \qquad\text{and }=\cot(k_h/2)\ \text{ on the half-zone of \cref{def:bethe-soft-limit}.}
\end{equation}
The remainders are uniform on compact hard sets: for $0<a<b<\pi$ there are
finite explicit constants $C_\delta(a,b)$, $C_S(a,b)$, of size $O(a^{-2})$ as
$a\downarrow0$ at fixed $b$, with
$\abs{R_\delta}\le C_\delta(a,b)\abs{k_s}^3$ and
$\abs{R_S}\le C_S(a,b)\abs{k_s}^3$ on the compactum
$\{\abs{k_s}\le\varepsilon_{ab},\ a\le\abs{k_h}\le b\}$.
\end{theorem}
\begin{scope}
The theorem is about the exact connected two-magnon multiplier of this model.
It does \emph{not} promote a process-independent soft theorem, and asserts
nothing about amplitudes generated by an arbitrary quasi-local hard process:
that would be the universality statement discussed in \cref{sec:universality}
and it is not established here. It is generalised --- not replaced --- by
\cref{thm:spin-s-slope}.
\end{scope}
\provenance{S2-2body}{\texttt{theory/soft-current-recon.md} steps (1.1)--(1.5)
together with \texttt{theory/oracle-bethe.md} step (1.3)}%
{\texttt{theory/checks/soft\_current\_recon\_check.py} (maximum operator and
form-factor residual $1.560\cdot10^{-14}$, quadratic-phase fit error
$2.167\cdot10^{-10}$), \texttt{oracle\_bethe\_check.py},
\texttt{ml2\_completeness\_check.py}, and the numerics scan
\texttt{fm-displacement-scan}}{\texttt{theory/verdicts/corpus-r2.md}
adjudication}

\emph{Idea of proof.} Write the complexified broken charge and its current,
\begin{equation}
  Q_k=\sum_xe^{ikx}S^-_x,\qquad
  j^-_{x,x+1}=-[h_{x,x+1},S^-_x]=\tfrac{J}{2}(S^-_{x+1}-S^-_x)P_{x,x+1},
  \qquad J^-_k=\sum_xe^{ikx}j^-_{x,x+1},
\end{equation}
so that $[H,Q_k]=(e^{ik}-1)J^-_k$. Acting with $Q_{k_s}$ on a hard magnon
produces a two-magnon state whose two ordered-chamber coefficients are
\emph{equal}; an on-shell scattering eigenstate has instead the ratio
$S_{12}$. The exact mismatch is a single branch term,
\begin{equation}
  Q_{k_s}\ket{k_h}_N-\ket{B^{\mathrm{in}}}=(1-S_{12})\ket{P_{12}},
  \label{eq:R8-mismatch}
\end{equation}
with $\ket{P_{12}}=N^{-1/2}\sum_{x<y}e^{i(k_sx+k_hy)}\ket{x,y}$. Now the soft
expansion is obtained without any closed formula: normalise the incoming
coefficient to one, call the unknown outgoing coefficient $s(k_s,k_h)$, and
impose only that the contact bond residual vanish. With $z_s=e^{ik_s}$ and
$z_h=e^{ik_h}$ this is one linear equation,
\begin{equation}
  (2z_h-z_sz_h-1)\,s+(2z_s-z_sz_h-1)=0 ,
  \label{eq:contact-equation}
\end{equation}
which for fixed $k_h$ has a unique solution analytic at $k_s=0$; expanding it
gives $s=1+2ik_s+[i\cot(k_h/2)-2]k_s^2+O(k_s^3)$, and taking the continuous
logarithm gives the phase. All hard-momentum dependence cancels from the
linear coefficient. Compactness of the hard set keeps $\abs{z_h-1}$ bounded
below, which is what makes Taylor's remainder uniform. Finally the current
identity supplies the mechanism behind the number $2$: at zero soft momentum
\begin{equation}
  \bra{k_h}Q_0^{\dagger}J^-_0\ket{k_h}_N=2iJ\sin k_h=2i\,v_h ,
  \label{eq:ward-residue}
\end{equation}
so a hard external-pole reduction, which divides by the linear energy shift
$v_hk_s$, cancels the hard velocity and leaves the coefficient $2$.

\begin{remark}[The number $2$ is a displacement, not a scattering length]
\label{rem:not-a-scattering-length}
The magnitude of the linear coefficient in \cref{eq:delta-phys} is the
two-site Wigner phase-slope, i.e.\ a spatial-displacement coefficient of the
soft packet, and its sign is the channel label $\sgn(v_h-v_s)$; the number is
not hard-momentum-independent in an invariant sense, and only the half-zone of
\cref{def:bethe-soft-limit} reduces it to $+2$. It is also not a scattering
length: in the joint soft limit $k_1=\varepsilon x$, $k_2=\varepsilon y$ with
$x\ne y$, the phase obeys $\delta_{12}/\varepsilon\to 2xy/(y-x)$, which
vanishes but is not a constant multiple of $k_1-k_2$, so no nonzero
relative-momentum scattering length exists. Two further boundary facts are
worth recording: as $k_h\downarrow0$ the quadratic coefficient diverges like
$2/k_h$, so truncation requires $\abs{k_s}\ll k_h$; and as $k_h\to\pi$ the
algebraic ratio stays regular but the kinematic ordering $v_h>v_s$ degenerates,
requiring $k_s\ll\pi-k_h$.
\end{remark}

\begin{remark}[Momentum is the good coordinate]
\label{rem:puiseux}
In the energy variable the same expansion is only Puiseux and remembers the
sign of the soft momentum through $\sigma=\sgn(k_s)$:
\begin{equation}
  \delta_{12}=2\sigma\sqrt{2\omega_s/J}
    +\frac{2v_h}{J\omega_h}\,\omega_s+O(\omega_s^{3/2}),
\end{equation}
with an explicit compact remainder bound. Both signed momentum limits give one,
but the energy coordinate forgets $\sigma$ and the connected term is
$O(\sqrt{\omega_s})$ rather than a Taylor series.
\end{remark}

\begin{refutedclaim}[The old label ``$S2$'' for the minimal two-body core]
\label{ref:s2-label}
\statusRefuted\ The identifier $S2$ was withdrawn --- purely as a
duplicate-status guard, so that one statement does not carry two rows --- and
its content is \cref{thm:two-body-soft}, which is the row that is proved.
\end{refutedclaim}
\provenance{S2 (superseded label; use S2-2body)}{---}{---}%
{superseded by the S2-2body row of \texttt{claims/CLAIMS.md}}

\subsection{The exact soft slope at every site spin}

The spin-$1/2$ result above turns out not to be an accident of $SU(2)$
spin-half algebra. The next theorem is the campaign's sharpest exact
statement.

\begin{theorem}[Exact two-magnon ratio and soft slope at every site spin]
\label{thm:spin-s-slope}
\statusProved\ Let $S\in\{1/2,1,3/2,\dots\}$, $J>0$, and take on the infinite
chain the fully polarised ferromagnet
\begin{equation}
  H_S=-J\sum_x\bigl(\bm{S}_x\cdot\bm{S}_{x+1}-S^2\bigr).
\end{equation}
Use the ordered chamber, Fourier sign, and coefficient ratio
$S_{12}=A_{12}/A_{21}$ of \cref{def:bethe-convention}. Put $z_j=e^{ik_j}$ and
\begin{equation}
  a:=1+z_1z_2,\qquad b:=z_1+z_2,\qquad \mu:=(2S-1)a+b .
\end{equation}
On the regular scattering domain --- real $k_1,k_2\in(-\pi,\pi]$ with
$k_1\ne k_2 \pmod{2\pi}$, $ab\ne0$, and $z_2\mu-Sab\ne0$ --- the two-magnon
eigenvalue equation \emph{by itself} fixes
\begin{equation}
  \boxed{\;S_{12}(k_1,k_2)=\frac{S\,ab-z_1\mu}{z_2\mu-S\,ab}\;}
  \label{eq:spin-s-ratio}
\end{equation}
with $\abs{S_{12}}=1$ and $S_{12}(k_2,k_1)=S_{12}(k_1,k_2)^{-1}$. Setting
$k_1=k_s$, $k_2=k_h$ with $0<\abs{k_h}<\pi$ fixed and signed $k_s\to0$, and
using the dispersion $\omega_S(k)=2JS(1-\cos k)$, $v_S(k)=2JS\sin k$ and the
physical channel rule of \cref{def:bethe-convention}, the physical phase
obeys, locally uniformly on compact hard sets on which the channel is fixed,
\begin{equation}
  \delta_{\mathrm{phys}}(k_s,k_h)=\frac{\sgn(v_h-v_s)}{S}\,k_s+O(k_s^2),
  \qquad
  S_{\mathrm{phys}}=1+i\,\frac{\sgn(v_h-v_s)}{S}\,k_s+O(k_s^2),
\end{equation}
so that
\begin{equation}
  \boxed{\;
   \left.\frac{\partial\delta_{\mathrm{phys}}}{\partial k_s}\right|_{k_s=0}
   =\frac{\sgn(v_h-v_s)}{S}\;}
  \label{eq:soft-slope-law}
\end{equation}
which on the half-zone of \cref{def:bethe-soft-limit} is exactly $1/S$. At
$S=1/2$ the ratio \cref{eq:spin-s-ratio} collapses to
\cref{eq:exact-ratio-half} and the slope to $2\sgn(v_h-v_s)$, recovering
\cref{thm:two-body-soft}.
\end{theorem}
\begin{scope}
No integrability hypothesis and no Bethe completeness are used: the two-body
wave is validated by direct substitution into every configuration class. What
is proved is the exact \emph{unit-charge two-body} slope, and only that. The
theorem does not prove endpoint ($k_h\to0,\pi$) or equal-velocity limits; it
does not prove Bethe completeness at spin $S$; it does not prove a
process-independent soft theorem; it does not prove the memory half of the
triangle (see \cref{sec:memory-index}); it says nothing about a hard
excitation of charge larger than one; and it therefore does not promote the
combined slope-and-memory conjecture, which remains a conjecture.
\end{scope}
\provenance{S2-2body-S}{\texttt{theory/spin-s-twomagnon.md}, steps
(1.1)--(1.5)}{\texttt{theory/checks/spin\_s\_slope\_check.py} (the \texttt{--red}
mutation replaces $2S-1$ by $2S+1$ in $\mu$ and must fail);
\texttt{numerics/results/spin1-bc-falsifier.json};
\texttt{numerics/results/spin1-bc-crosscheck.json}}%
{\texttt{theory/verdicts/spin-s-r1.md} (PASS)}

\emph{Idea of proof.} In occupation variables $n_x=S-S^z_x$ the Hamiltonian is
$H_S=2JS\sum_xn_x-J\sum_xn_xn_{x+1}-\tfrac{J}{2}\sum_x(S^+_xS^-_{x+1}+S^-_xS^+_{x+1})$.
For $S\ge1$ the two-magnon sector has two independent amplitudes: the
separated wave $\psi(x,y)$ on $x<y$ and a genuine double-occupancy amplitude
$d(x)$. There are therefore three exact coordinate equations --- separated,
adjacent, and doubly occupied --- and, with $g=\sqrt{S(2S-1)}$, they read
\begin{align}
  E\psi(x,y)&=4JS\psi(x,y)-JS\bigl[\psi(x-1,y)+\psi(x+1,y)+\psi(x,y-1)+\psi(x,y+1)\bigr],\\
  E\psi(x,x+1)&=J(4S-1)\psi(x,x+1)-JS\bigl[\psi(x-1,x+1)+\psi(x,x+2)\bigr]
    -Jg\bigl[d(x)+d(x+1)\bigr],\\
  E\,d(x)&=4JS\,d(x)-Jg\bigl[\psi(x-1,x)+\psi(x,x+1)\bigr].
\end{align}
One takes the unsymmetrised two-plane-wave extension
$\Psi(x,y)=Az_1^xz_2^y+Bz_2^xz_1^y$ on all of $\Z^2$, together with
$d(x)=\rho(z_1z_2)^x$. The separated equation fixes the energy additively; the
double-occupancy equation determines $\rho$ in terms of the amplitudes; and
eliminating $\rho$ from the adjacent equation leaves the single contact
relation $W\mu=S\Sigma ab$ with $\Sigma=A+B$ and $W=Az_2+Bz_1$, which is
exactly \cref{eq:spin-s-ratio}. Unitarity follows by writing the ratio in
centre/relative variables, where it takes the form $n/(-\bar n)$.
Differentiating $\log S_{12}$ at $z_1=1$ gives $i/S$, whence the slope.

One trap deserves recording, because it is what makes the answer $1/S$ rather
than $2$: the free extension $\Psi$ is \emph{not} symmetric across the
diagonal unless $A=B$, so one may not replace $\Psi(x+1,x)$ by $\Psi(x,x+1)$
when extending the separated equation to coincident arguments. That illicit
replacement collapses the two contact conditions into the spin-half one and
falsely yields slope $2$ at every $S$. The valid argument keeps the extension
unsymmetrised, introduces $d(x)$ as an independent physical amplitude, and
imposes both contact equations.

\subsection{Consistency with the oracle, and the pending numerical test}

\begin{proposition}[The reconstruction agrees with the passed oracle]
\label{thm:oracle-consistency}
\statusProved\ On the half-zone of \cref{def:bethe-soft-limit}, the expansion
obtained from the local current and the single contact equation
\cref{eq:contact-equation} agrees term by term --- including the sign
convention and the value $2$ of the linear coefficient there --- with the
independently derived closed-form expansions \cref{eq:delta-phys,eq:S-phys} of
the Bethe oracle.
\end{proposition}
\begin{scope}
This establishes the equality of two formulas, obtained by two independent
routes. It does not establish process universality: agreement of the
contact-equation route with the closed Bethe route says nothing about
amplitudes generated by other hard processes.
\end{scope}
\provenance{OR1}{\texttt{theory/soft-current-recon.md} step
(1.5).(2).1}{\texttt{theory/checks/oracle\_bethe\_check.py},
\texttt{theory/checks/soft\_current\_recon\_check.py}}%
{\texttt{theory/verdicts/corpus-r2.md}}

\begin{conjecture}[Excitation-ansatz amplitudes reproduce the Bethe phase]
\label{conj:excitation-ansatz-bethe}
\statusConjecture\ In a variational matrix-product excitation-ansatz
computation of two-magnon scattering in the model of
\cref{def:fm-oracle-model}, the extracted magnon amplitudes reproduce the
exact Bethe ratio $S(k)$ as $k\to0$.
\end{conjecture}
\begin{scope}
This numerical test has not been run. It is stated so that the tensor-network
side of the campaign (\cref{sec:numerics}) has an unambiguous, falsifiable
target inherited from \cref{thm:two-body-soft}; no evidence for it is claimed
here.
\end{scope}
\provenance{N1}{---}{--- (not yet run)}{---}
